\section*{Chapter \ref{ch:pl:reference-data-and-symbology-services} --- Reference Data and Symbology Services}

\begin{solution}{exo:pl:reference-data-and-symbology-services:1}
$384 / 10\,493 = 3.7\%$. It happens on the dates when the ticker belonged to another listing or to none: before a rename to it,
after a rename away from it, and before a reuse of it (\cref{prop:pl:reference-data-and-symbology-services:today}).
\end{solution}

\begin{solution}{exo:pl:reference-data-and-symbology-services:2}
The CUSIP stayed with the security through the ticker change, so a history keyed by it continues unbroken where one keyed by
ticker breaks (or joins another issuer's). It is still an external identifier: it changes with some corporate actions, is issued
by one numbering agency for one market, and does not cover every instrument the firm trades; the firm keys by its own permanent
identifier and stores the CUSIP, ISIN, FIGI and tickers as dated attributes.
\end{solution}

\begin{solution}{exo:pl:reference-data-and-symbology-services:3}
1 on day 385 (nothing known), 1 on day 386 (announced but not confirmed), 4 from day 389 (confirmed), and 4 on day 412.
\end{solution}

\begin{solution}{exo:pl:reference-data-and-symbology-services:4}
The mistyped lot, 10, on all six days: precedence copies the preferred source's errors. The conflict report catches it, since
vendor B says 100 on each of those days. A validation rule would have stopped it: a lot size that changes by a factor of ten
without a corporate action, or that is not a round lot of the listing's market, is quarantined for review and the \omterm{def:pl:reference-data-and-symbology-services:golden}{golden copy} keeps
the previous value meanwhile.
\end{solution}

\begin{solution}{exo:pl:reference-data-and-symbology-services:5}
As known on day 255: the fund itself (the acquisition was not yet known). As known on day 262: the holding company, through the
bank. The reproduced report must show the first answer, what the firm knew when it printed the report; the second is a correction
to report separately.
\end{solution}

\begin{solution}{exo:pl:reference-data-and-symbology-services:6}
Unadjusted splits turn each split into a fall of $1 - 1/r$ in the price series (a four-for-one looks like $-75\%$): on the right
listings the basket returns $-22.3\%$ instead of $+23.5\%$, 45.8 points. Today's mapping replaces five intended listings by the
listings that now carry their tickers, with their own returns from their own first day: split-adjusted, $+29.2\%$, 5.7 points
above the truth. Together they give $-16.6\%$, 40.1 points below: the effects do not add because the unadjusted splits fall on
the intended listings, some of which today's mapping has replaced by listings with no split.
\end{solution}

\begin{solution}{exo:pl:reference-data-and-symbology-services:7}
150 instrument-days. Vendor B is wrong on 153, but on three of them it has no ticker at all for the new listing, and the \omterm{def:pl:reference-data-and-symbology-services:golden}{golden copy}
falls back to vendor A, which is right.
\end{solution}

\begin{solution}{exo:pl:reference-data-and-symbology-services:8}
Returns computed from adjusted prices are fine, but a signal that uses price levels is not: with today's factors, the price
before a split is divided by a ratio that was not known then (not even announced), so a rule such as ``price below \$5'' or
``near a round number'' sees prices the market never showed. Adjust with the factors known at each decision time, or use
returns only.
\end{solution}

\begin{solution}{pb:pl:reference-data-and-symbology-services:1}
\begin{enumerate}
\item 12 renames, 5 reuses, 20 splits, one of them cancelled.
\item The valid date (from when the fact is true) and the knowledge time (when the firm learned it): the first to ask about the
past, the second to reproduce what was believed then.
\item Two renames two days late, and one rename not until three days after the old ticker was reused: 146 days of a wrong
ticker for one listing.
\item A lot size of 10 instead of 100 for six days.
\item The vendors' values are the evidence for resolving conflicts, for measuring each source, and for rebuilding the \omterm{def:pl:reference-data-and-symbology-services:golden}{golden copy}
under a new rule.
\item Vendor A never, vendor B on 153 instrument-days, the \omterm{def:pl:reference-data-and-symbology-services:golden}{golden copy} never.
\item Vendor A is first for lots and was wrong; the \omterm{def:pl:reference-data-and-symbology-services:golden}{golden copy} takes its value.
\item 150 for tickers, 6 for lots.
\item The listing that renamed on day 33 (vendor A: its new ticker; vendor B: still the old one), on the day its old ticker went to
a new listing.
\item Keep A first for tickers (it was never wrong on them), and add validation for lots rather than moving the precedence: A is
right about lots everywhere but on those six days.
\item 3.7\%.
\item Listing 56 until day 32, no listing from day 33 to day 175, listing 100 from day 176; vendor B keeps it on listing 56 until
day 178.
\item 1 before the confirmation on day 389, 4 after.
\item It never enters the factor; the announcement and the cancellation stay in its history.
\item The fund itself, as known on day 255; the holding company, as known on day 262.
\item \textbf{Named result.} The 20-ticker basket returns $+23.5\%$ point in time, equal to the truth; with today's mapping and
unadjusted prices it returns $-16.6\%$, 40.1 points wrong, of which unadjusted splits alone account for $-45.8$ points and today's
mapping alone for $+5.7$.
\item The splits: a $-75\%$ day is visible at once. Today's mapping is the more dangerous: the wrong listings' prices are real and
plausible.
\item The listing of each symbol at each decision date and the adjustment factors, both as known at the decision time.
\item Resolve it on the first day from an exchange notice, record the resolution, and raise vendor B's performance on renames with
the vendor.
\item Because the firm must answer both what was true on a date and what it knew when it acted, and the two differ whenever a fact
arrives late or is corrected.
\end{enumerate}
\end{solution}

\begin{solution}{iq:pl:reference-data-and-symbology-services:1}
Tickers change and are reused: a history keyed by ticker splits one company's past across symbols and joins different companies
under one symbol. Key by a permanent identifier and keep tickers as dated attributes.
\iqlookfor{Reuse and renames; a permanent key.}
\end{solution}

\begin{solution}{iq:pl:reference-data-and-symbology-services:2}
By a stated precedence for that field, measured on each vendor's accuracy, plus validation; report the conflict to operations and
resolve it from evidence. Keep both vendors' values, the winner's provenance and any override, bitemporally.
\iqlookfor{Precedence with provenance, conflicts reported, nothing discarded.}
\end{solution}

\begin{solution}{iq:pl:reference-data-and-symbology-services:3}
Adjusting past prices with factors known only later (a split announced after the decision) changes price levels the strategy saw,
and a corporate action applied from its effective date though learned later puts knowledge in the past. Store actions as versioned
events and ask for factors as known at each decision.
\iqlookfor{Knowledge time of the adjustment.}
\end{solution}

\begin{solution}{iq:pl:reference-data-and-symbology-services:4}
Limits are set on ultimate parents. If a counterparty is acquired, the exposure aggregates to a new parent from the acquisition
date, but the firm learned it later: a limit check made in between was right with what was known, and must be reproducible that
way, while the correction shows the exposure the firm actually had.
\iqlookfor{Valid against knowledge time, applied to an aggregation.}
\end{solution}

\begin{solution}{iq:pl:reference-data-and-symbology-services:5}
Sources: exchanges' reference files, two or more vendors, a corporate-action feed, a legal-entity source. Permanent internal
identifiers with every external one (tickers per venue, ISIN, CUSIP, FIGI, vendor codes) as dated attributes; bitemporal storage;
a \omterm{def:pl:reference-data-and-symbology-services:golden}{golden copy} per field with precedence, validation, overrides and provenance; a conflict queue; corporate actions as versioned
events. Consumers query a service by identifier or symbol with a valid date and a knowledge time, never a copied file.
\iqlookfor{Identifiers, bitemporality, \omterm{def:pl:reference-data-and-symbology-services:golden}{golden copy} with provenance, a query interface.}
\end{solution}
